% =======================================================
% 2. ANALISI DEI RISCHI
% =======================================================
\section{Analisi dei Rischi}
\subsection{Nomenclatura e lettura della tabella}
Ogni rischio è identificato da un codice composto da due lettere indicanti la categoria e un numero progressivo:
\begin{itemize}
	\item \textbf{RR}: Rischi legati ai Requisiti (ambiguità, scope creep, qualità, scadenze)
	\item \textbf{RT}: Rischi Tecnici (tecnologie, integrazione, strumenti, debito tecnico)
	\item \textbf{RO}: Rischi Organizzativi (disponibilità, comunicazione, carico, motivazione, comprensione)
\end{itemize}

Per ciascun rischio vengono indicati:
\begin{itemize}
    \item \textbf{Codice:}  Identificativo univoco del rischio (es. RR1, RT2, RO5).
    \item \textbf{Descrizione:}  Descrizione sintetica del rischio.
    \item \textbf{Probabilità di occorrenza:}  Valutazione della probabilità che il rischio si manifesti (Alta / Media / Bassa).
    \item \textbf{Impatto:}  Valutazione dell'impatto che il rischio avrebbe sul progetto (Alta / Media / Bassa).
	\item \textbf{Piano di mitigazione:} Azioni preventive volte a ridurre la probabilità o l'impatto del rischio.
	\item \textbf{Piano di contingenza:} Azioni da intraprendere nel caso in cui il rischio si verifichi.
\end{itemize}

\subsection{RR - Rischi Legati ai Requisiti}

\begin{table}[H]
	\centering
	\renewcommand{\arraystretch}{1.3}
	\begin{tabularx}{\textwidth}{|c|X|}
		\hline
		\textbf{Codice} & \textbf{RR1} \\
		\hline
		\textbf{Descrizione} & Requisiti ambigui o cambiamenti frequenti da parte del committente \\
		\hline
		\textbf{Probabilità / Impatto} & Media/Alta \\
		\hline
		\textbf{Piano di Mitigazione} & Versionare i requisiti nel VCS; approvare formalmente le baseline; mantenere confronto continuo con il proponente \\
		\hline
		\textbf{Piano di Contingenza} & Attivazione del processo di change request; revisione backlog; rinegoziazione delle priorità e della pianificazione \\
		\hline
	\end{tabularx}
	\caption{Scheda rischio RR1}
\end{table}

\begin{table}[H]
	\centering
	\renewcommand{\arraystretch}{1.3}
	\begin{tabularx}{\textwidth}{|c|X|}
		\hline
		\textbf{Codice} & \textbf{RR2} \\
		\hline
		\textbf{Descrizione} & Scope creep: aggiunta incontrollata di funzionalità \\
		\hline
		\textbf{Probabilità / Impatto} & Alta/Alta \\
		\hline
		\textbf{Piano di Mitigazione} & Definizione chiara dello scope RTB; backlog prioritizzato; approvazione preventiva delle nuove feature \\
		\hline
		\textbf{Piano di Contingenza} & Posticipo delle funzionalità non essenziali alla PB; riduzione dello scope MVP \\		
		\hline
	\end{tabularx}
	\caption{Scheda rischio RR2}
\end{table}

\begin{table}[H]
	\centering
	\renewcommand{\arraystretch}{1.3}
	\begin{tabularx}{\textwidth}{|c|X|}
		\hline
		\textbf{Codice} & \textbf{RR3} \\
		\hline
		\textbf{Descrizione} & Qualità insufficiente del prodotto (bug, mancanza di test) \\
		\hline
		\textbf{Probabilità / Impatto} & Media/Alta \\
		\hline
		\textbf{Piano di Mitigazione} & Definition of Done con criteri di qualità espliciti; code review sistematiche; copertura minima di test definita \\
		\hline
		\textbf{Piano di Contingenza} & Sprint dedicato al bug fixing; riallocazione del tempo verso attività di verifica e validazione \\
		\hline
	\end{tabularx}
	\caption{Scheda rischio RR3}
\end{table}

\begin{table}[H]
	\centering
	\renewcommand{\arraystretch}{1.3}
	\begin{tabularx}{\textwidth}{|c|X|}
		\hline
		\textbf{Codice} & \textbf{RR4} \\
		\hline
		\textbf{Descrizione} & Mancato rispetto delle scadenze (RTB o PB) \\
		\hline
		\textbf{Probabilità / Impatto} & Media/Alta \\
		\hline
		\textbf{Piano di Mitigazione} & Pianificazione conservativa con buffer temporali; monitoraggio settimanale dell'avanzamento e della burndown chart \\
		\hline
		\textbf{Piano di Contingenza} & Riduzione dello scope MVP; rinegoziazione delle milestone interne; riallocazione delle risorse \\
		\hline
	\end{tabularx}
	\caption{Scheda rischio RR4}
\end{table}

\begin{table}[H]
	\centering
	\renewcommand{\arraystretch}{1.3}
	\begin{tabularx}{\textwidth}{|c|X|}
		\hline
		\textbf{Codice} & \textbf{RR5} \\
		\hline
		\textbf{Descrizione} & Ritardo nella validazione dei requisiti da parte del proponente \\
		\hline
		\textbf{Probabilità / Impatto} & Media/Alta \\
		\hline
		\textbf{Piano di Mitigazione} & Organizzare incontri periodici di validazione incrementale; condividere frequentemente gli avanzamenti \\
		\hline
		\textbf{Piano di Contingenza} & Congelamento temporaneo dei requisiti già condivisi; prosecuzione dello sviluppo sulle funzionalità validate \\
		\hline
	\end{tabularx}
	\caption{Scheda rischio RR5}
\end{table}

\subsection{RT -Rischi Tecnici}
\begin{table}[H]
	\centering
	\renewcommand{\arraystretch}{1.3}
	\begin{tabularx}{\textwidth}{|c|X|}
		\hline
		\textbf{Codice} & \textbf{RT1} \\
		\hline
		\textbf{Descrizione} & Inesperienza con le tecnologie scelte (framework, librerie, API) \\
		\hline
		\textbf{Probabilità / Impatto} & Media/Alta \\
		\hline
		\textbf{Piano di Mitigazione} & Formazione iniziale, spike tecnici e studio delle tecnologie candidate; pair programming \\
		\hline
		\textbf{Piano di Contingenza} & Riduzione dello scope tecnico; sostituzione di librerie o componenti troppo complessi \\
		\hline
	\end{tabularx}
	\caption{Scheda rischio RT1}
\end{table}

\begin{table}[H]
	\centering
	\renewcommand{\arraystretch}{1.3}
	\begin{tabularx}{\textwidth}{|c|X|}
		\hline
		\textbf{Codice} & \textbf{RT2} \\
		\hline
		\textbf{Descrizione} & Problemi di integrazione tra componenti (frontend/backend, API esterne) \\
		\hline
		\textbf{Probabilità / Impatto} & Media/Media \\
		\hline
		\textbf{Piano di Mitigazione} & Definizione preventiva delle interfacce e dei contratti API; utilizzo di mock e test di integrazione \\
		\hline
		\textbf{Piano di Contingenza} & Sessioni straordinarie di integrazione; revisione dell'architettura dei componenti coinvolti \\
		\hline
	\end{tabularx}
	\caption{Scheda rischio RT2}
\end{table}

\begin{table}[H]
	\centering
	\renewcommand{\arraystretch}{1.3}
	\begin{tabularx}{\textwidth}{|c|X|}
		\hline
		\textbf{Codice} & \textbf{RT3} \\
		\hline
		\textbf{Descrizione} & Indisponibilità o instabilità di strumenti/servizi esterni (GitHub, CI/CD, cloud) \\
		\hline
		\textbf{Probabilità / Impatto} & Bassa/Media \\
		\hline
		\textbf{Piano di Mitigazione} & Predisporre strumenti alternativi e backup locali; documentare workflow indipendenti dai servizi esterni \\
		\hline
		\textbf{Piano di Contingenza} & Passaggio temporaneo ad ambienti locali; aggiornamento della pianificazione e sospensione delle pipeline dipendenti \\
		\hline
	\end{tabularx}
	\caption{Scheda rischio RT3}
\end{table}

\begin{table}[H]
	\centering
	\renewcommand{\arraystretch}{1.3}
	\begin{tabularx}{\textwidth}{|c|X|}
		\hline
		\textbf{Codice} & \textbf{RT4} \\
		\hline
		\textbf{Descrizione} & Debito tecnico accumulato che rallenta gli sprint successivi \\
		\hline
		\textbf{Probabilità / Impatto} & Media/Media \\
		\hline
		\textbf{Piano di Mitigazione} & Refactoring incrementale continuo; code review; riserva di capacità dedicata al debito tecnico \\
		\hline
		\textbf{Piano di Contingenza} & Sprint dedicato alla manutenzione correttiva e al consolidamento tecnico \\
		\hline
	\end{tabularx}
	\caption{Scheda rischio RT4}
\end{table}

\begin{table}[H]
	\centering
	\renewcommand{\arraystretch}{1.3}
	\begin{tabularx}{\textwidth}{|c|X|}
		\hline
		\textbf{Codice} & \textbf{RT5} \\
		\hline
		\textbf{Descrizione} & Le tecnologie selezionate potrebbero non soddisfare pienamente i requisiti del progetto \\
		\hline
		\textbf{Probabilità / Impatto} & Media/Alta \\
		\hline
		\textbf{Piano di Mitigazione} & Analisi comparativa delle tecnologie; sviluppo del PoC; confronto continuo con il proponente \\
		\hline
		\textbf{Piano di Contingenza} & Sostituzione delle componenti critiche; riduzione dello scope tecnico non essenziale \\
		\hline
	\end{tabularx}
	\caption{Scheda rischio RT5}
\end{table}

\subsection{RO - Rischi Organizzativi}

\begin{table}[H]
	\centering
	\renewcommand{\arraystretch}{1.3}
	\begin{tabularx}{\textwidth}{|c|X|}
		\hline
		\textbf{Codice} & \textbf{RO1} \\
		\hline
		\textbf{Descrizione} & Membro del team non disponibile per malattia o impegni accademici \\
		\hline
		\textbf{Probabilità / Impatto} & Media/Alta \\
		\hline
		\textbf{Piano di Mitigazione} & Condivisione continua delle conoscenze; documentazione costante delle attività; evitare single point of failure \\
		\hline
		\textbf{Piano di Contingenza} & Ridistribuzione del carico; riduzione dello scope dello sprint; aggiornamento della pianificazione \\
		\hline
	\end{tabularx}
	\caption{Scheda rischio RO1}
\end{table}

\begin{table}[H]
	\centering
	\renewcommand{\arraystretch}{1.3}
	\begin{tabularx}{\textwidth}{|c|X|}
		\hline
		\textbf{Codice} & \textbf{RO2} \\
		\hline
		\textbf{Descrizione} & Conflitti interni al team o scarsa comunicazione \\
		\hline
		\textbf{Probabilità / Impatto} & Bassa/Alta \\
		\hline
		\textbf{Piano di Mitigazione} & Daily standup regolari; retrospettive periodiche; definizione di convenzioni condivise \\
		\hline
		\textbf{Piano di Contingenza} & Mediazione del responsabile; eventuale coinvolgimento del docente o del proponente \\
		\hline
	\end{tabularx}
	\caption{Scheda rischio RO2}
\end{table}

\begin{table}[H]
	\centering
	\renewcommand{\arraystretch}{1.3}
	\begin{tabularx}{\textwidth}{|c|X|}
		\hline
		\textbf{Codice} & \textbf{RO3} \\
		\hline
		\textbf{Descrizione} & Distribuzione sbilanciata del carico di lavoro \\
		\hline
		\textbf{Probabilità / Impatto} & Media/Media \\
		\hline
		\textbf{Piano di Mitigazione} & Monitoraggio continuo delle ore e della velocity; revisione periodica dell'assegnazione task \\
		\hline
		\textbf{Piano di Contingenza} & Ribilanciamento immediato del backlog; redistribuzione delle responsabilità \\
		\hline
	\end{tabularx}
	\caption{Scheda rischio RO3}
\end{table}

\begin{table}[H]
	\centering
	\renewcommand{\arraystretch}{1.3}
	\begin{tabularx}{\textwidth}{|c|X|}
		\hline
		\textbf{Codice} & \textbf{RO4} \\
		\hline
		\textbf{Descrizione} & Perdita di motivazione o disimpegno di uno o più membri \\
		\hline
		\textbf{Probabilità / Impatto} & Bassa/Media \\
		\hline
		\textbf{Piano di Mitigazione} & Obiettivi incrementali chiari; monitoraggio del carico di lavoro; coinvolgimento continuo del team \\
		\hline
		\textbf{Piano di Contingenza} & Colloqui individuali; riallocazione task; supporto organizzativo del gruppo \\
		\hline
	\end{tabularx}
	\caption{Scheda rischio RO4}
\end{table}

\begin{table}[H]
	\centering
	\renewcommand{\arraystretch}{1.3}
	\begin{tabularx}{\textwidth}{|c|X|}
		\hline
		\textbf{Codice} & \textbf{RO5} \\
		\hline
		\textbf{Descrizione} & Scarsa comprensione dei requisiti o della visione di progetto \\
		\hline
		\textbf{Probabilità / Impatto} & Media/Alta \\
		\hline
		\textbf{Piano di Mitigazione} & Workshop iniziali di analisi; aggiornamento continuo del glossario e dei casi d'uso; confronto periodico con il proponente \\
		\hline
		\textbf{Piano di Contingenza} & Revisione straordinaria dei requisiti; riallineamento del backlog e delle priorità progettuali \\
		\hline
	\end{tabularx}
	\caption{Scheda rischio RO5}
\end{table}

\begin{table}[H]
	\centering
	\renewcommand{\arraystretch}{1.3}
	\begin{tabularx}{\textwidth}{|c|X|}
		\hline
		\textbf{Codice} & \textbf{RO6} \\
		\hline
		\textbf{Descrizione} & Sottostima del carico documentale e organizzativo associato al ruolo di Amministratore \\
		\hline
		\textbf{Probabilità / Impatto} & Alta/Media \\
		\hline
		\textbf{Piano di Mitigazione} & Distribuzione delle attività redazionali; revisione periodica della pianificazione delle ore \\
		\hline
		\textbf{Piano di Contingenza} & Aggiornamento del Preventivo a Finire; riallocazione del budget orario tra i membri del gruppo \\
		\hline
	\end{tabularx}
	\caption{Scheda rischio RO6}
\end{table}

\subsection{Matrice dei rischi}

\begin{table}[H]
\centering
\renewcommand{\arraystretch}{1.5}

\begin{tabular}{|c|c|c|c|}
\hline
\diagbox{Impatto}{Probabilità} & Bassa & Media & Alta \\
\hline

Alta
& --
& RR1, RR4, RR5, RT1, RT5, RO1, RO5
& RR2 \\
\hline

Media
& RT3, RO4
& RT2, RT4, RO3
& RO6 \\
\hline

Bassa
& --
& --
& -- \\
\hline

\end{tabular}

\caption{Matrice qualitativa dei rischi}
\end{table}

\subsection{Monitoraggio dei rischi}
I rischi vengono rivalutati al termine di ogni Sprint Review.
Il Responsabile monitora l'evoluzione dei rischi identificati,
verifica l'efficacia delle azioni di mitigazione adottate e aggiorna
il registro dei rischi sulla base di eventuali nuovi elementi emersi
durante lo sviluppo del progetto.